(1) 设模式场$\psi(r,\phi,z)=F(r)e^{il\phi}e^{i\beta z}$，每个均匀区满足$F''+r^{-1}F'+[n_j^2k_0^2-\beta^2-l^2/r^2]F=0$。导引模要求$n_2k_0<\beta<n_1k_0$，芯内有限解为$F_1=A J_l(ur/a)$，包层指数衰减解为$F_2=B K_l(wr/a)$，其中$u^2=a^2(n_1^2k_0^2-\beta^2)>0$、$w^2=a^2(\beta^2-n_2^2k_0^2)>0$。
(2) 在$r=a$连续标量场和径向导数，分别得到$AJ_l(u)=BK_l(w)$、$A uJ_l'(u)=B w K_l'(w)$。消元得$uJ_l'(u)/J_l(u)=wK_l'(w)/K_l(w)$。用Bessel导数递推可写成用于LP谱求根的超越方程；完整矢量Maxwell边界还存在偏振相关的细分模式，弱导引时才近似简并。
(3) 两定义平方相加，$V^2=u^2+w^2=a^2k_0^2(n_1^2-n_2^2)$，即$V=(2\pi a/\lambda)\sqrt{n_1^2-n_2^2}$。LP$_{01}$基模无有限$V$截止，下一组LP$_{11}$在$w\to0$时的边界条件简化为$J_0(u)=0$，第一个正零点$u=2.404826\ldots$，因此严格单模近似条件为$V<2.405$。
(4) 定义归一化传播常数$b=(\beta^2-n_2^2k_0^2)/(n_1^2k_0^2-n_2^2k_0^2)=w^2/V^2$，导引模有$0<b<1$；趋截止时$w\to0$、$b\to0$，模式向包层延伸，衰减长度$a/w$发散。大数值孔径、强折射率对比或显著材料色散时需解全矢量混合HE/EH模式，不能机械使用LP单模阈值。
